\title{Group Sequence Policy Optimization}

\begin{abstract}
We present Group Sequence Policy Optimization (GSPO), a reinforcement learning algorithm specifically designed to offer stability, computational efficiency, and high performance for training large language models. Departing from approaches that utilize token-level importance ratios, GSPO introduces a sequence-level perspective—computing importance ratios from the probability of entire sequences, and applying sequence-based clipping, reward assignment, and optimization. Our experiments indicate that GSPO outperforms the GRPO algorithm, enhancing both training efficiency and overall performance, while also ensuring greater stability in Mixture-of-Experts (MoE) RL settings. Furthermore, GSPO's methodology opens opportunities for simplifying RL system infrastructure. These advantages have played a critical role in the advancements observed in the most recent Qwen3 models.
\end{abstract}

\section{Introduction}
\label{sec:intro}

Reinforcement learning (RL) has become a foundational framework for advancing the capabilities of language models at scale \citep{o1,dpsk-r1,qwq32b,qwen3}. Leveraging large-scale RL enables these models to approach complex tasks—such as those seen in advanced mathematics and programming—by engaging in more intricate and prolonged reasoning.

Achieving successful scaling of RL with increased computational resources critically depends on the ability to ensure stable and resilient training trajectories. Nevertheless, the prevailing RL methods at the forefront of the field, including GRPO \citep{grpo}, encounter pronounced instability when applied to the training of large-scale language models. This often culminates in sudden and irreparable model degradation, as noted in prior studies \citep{qwen3, minimax-m1}. Such fragility fundamentally constrains the progress toward enhancing language model performance through sustained RL-based optimization.

In this work, we demonstrate that the principal source of instability in GRPO arises from a fundamental misapplication and eventual invalidity of importance sampling weights within its architectural framework. This flaw introduces substantial training noise with high variance, which accrues as the response lengthens and is exacerbated by the use of a clipping technique, ultimately leading to the collapse of the model.

To overcome these fundamental issues, we introduce \textbf{Group Sequence Policy Optimization (GSPO)}, a novel reinforcement learning algorithm tailored for large language model training. GSPO’s principal contribution is a theoretically justified formulation of the importance ratio grounded in sequence likelihood \citep{click}, thereby preserving the core integrity of importance sampling. Furthermore, GSPO leverages normalized rewards, computed as advantages across several responses to the same query, facilitating coherent optimization at the sequence level in line with reward assignment.

Our experimental results highlight that GSPO substantially surpasses GRPO regarding stability, training efficiency, and overall effectiveness. Notably, GSPO intrinsically addresses the instability issues commonly encountered in reinforcement learning for large Mixture-of-Experts (MoE) architectures, thus obviating the reliance on intricate stabilization techniques and indicating possibilities for streamlined RL frameworks. These strengths of GSPO have played a pivotal role in driving the notable performance gains observed in the most recent Qwen3 models. Looking ahead, we anticipate that GSPO will serve as a stable and scalable foundational method, supporting further progress in large-scale reinforcement learning with language models.

\section{Preliminaries}

\paragraph{Notation}
Throughout this work, we refer to an autoregressive language model with parameters $\theta$ as a policy $\pi_\theta$. Here, $x$ represents a query, while $\mathcal{D}$ denotes the set of all queries considered. For any response $y$ corresponding to a query $x$, the policy assigns the probability $\pi_\theta (y | x) = \prod_{t=1}^{|y|} \pi_\theta (y_t | x, y_{<t} )$, where $|y|$ indicates the length of $y$ in terms of tokens. We evaluate each query-response pair $(x, y)$ via a verifier $r$, which produces a reward $r(x, y) \in [ 0, 1 ]$.

\paragraph{Proximal Policy Optimization (PPO)} Using samples generated from the old policy $\pi_{\theta_\text{old}}$, PPO \citep{ppo} constrains the policy update within a proximal region of the old policy through the clipping mechanism. Specifically, PPO employs the following objective for policy optimization (we omit the KL regularization term hereinafter for brevity, as it is not the focus of this paper):

\begin{align}
\mathcal{J}_\text{PPO}(\theta) = \mathbb{E}_{ x \sim \mathcal{D},\, y \sim \pi_{\theta_\text{old}}( \cdot | x) }
\left[ \frac{1}{|y|} \sum_{t=1}^{|y|} 
\min \left( w_{t}(\theta) \widehat{A}_{t},  \, \mathrm{clip} \left( w_{t}(\theta), 1 - {\varepsilon}, 1 + {\varepsilon}\right) \widehat{A}_{t} \right)
\right],
\end{align}

In this context, the token importance ratio is given by $
w_{t}(\theta) = \frac{ \pi_{\theta} (y_{t} | x, y_{<t}) }{ \pi_{\theta_\text{old}} (y_{t} | x,y_{<t})} $, where the advantage term $\widehat{A}_{t}$ associated with $y_t$ is produced by a separate value model, and $\varepsilon$ denotes the clipping threshold for these importance ratios.

A fundamental issue with PPO in practical scenarios is its significant dependence on the value model. Typically, the value model is constructed to be nearly as large as the policy model, thereby substantially increasing both memory requirements and computational complexity. Moreover, PPO’s practical efficacy is closely tied to the accuracy of its value estimates. Crafting a dependable value model is an intrinsically difficult problem, and making such a model robust enough for extended outputs and more intricate tasks constitutes an even steeper challenge.

\paragraph{Group Relative Policy Optimization (GRPO)} GRPO \citep{grpo} bypasses the need for the value model by computing the relative advantage of each response within a group of responses to the same query. Specifically, GRPO optimizes the following objective:

\begin{align}
\label{equ:grpo}
\mathcal{J}_\text{GRPO}(\theta) = \mathbb{E}_{ x \sim \mathcal{D},\, \{y_i\}_{i=1}^G \sim \pi_{\theta_\text{old}}( \cdot | x) }
\left[ \frac{1}{G} \sum_{i=1}^{G} \frac{1}{|y_i|} \sum_{t=1}^{|y_i|} 
\min \left( w_{i,t}(\theta) \widehat{A}_{i,t},  \, \mathrm{clip} \left( w_{i,t}(\theta), 1 - {\varepsilon}, 1 + {\varepsilon}\right) \widehat{A}_{i,t} \right)
\right],
\end{align}

where $G$ is the number of generated responses to each query $x$ (i.e., the group size), and the importance ratio $w_{i,t}(\theta)$ and advantage $\widehat{A}_{i,t}$ of token $y_{i,t}$ are:

\begin{align}
    w_{i,t}(\theta)=\frac{ \pi_{\theta} (y_{i,t} | x, y_{i,<t}) }{ \pi_{\theta_\text{old}} (y_{i,t} | x,y_{i,<t})},\quad\ 
    \widehat{A}_{i,t} = \widehat{A}_{i} = \frac{r(x, y_i) - \mathrm{mean} \left( \{ r(x, y_i) \}_{i=1}^G \right) }{ \mathrm{std} \left( \{ r(x, y_i) \}_{i=1}^G \right) },
\end{align}

respectively, where all the tokens in $y_i$ share the same advantage as $\widehat{A}_{i}$.

\section{Motivation}
\label{sec:motivation}

As model architectures scale up, incorporating features like increased parameter counts, sparsity patterns typical in Mixture-of-Experts models, and longer generated outputs, reinforcement learning methods must employ large rollout batch sizes to ensure efficient use of computational hardware. To increase the number of learning updates per collected sample, standard techniques divide this substantial batch of rollout data into several smaller mini-batches for repeated gradient computations. This strategy, however, leads to an inherent off-policy learning scenario: the sequences $y$ utilized in gradient updates are drawn from a previous policy $\pi_{\theta_\text{old}}$, not the currently updated policy $\pi_\theta$. Consequently, clipping mechanisms such as those found in PPO and GRPO are essential; these restrict the influence of samples that are too ``off-policy'', thereby stabilizing gradient estimates.

While mechanisms like clipping aim to manage this off-policy discrepancy, we identify a more fundamental issue in GRPO: \textit{its objective is ill-posed}. This problem becomes particularly acute when training large models on long-response tasks, leading to catastrophic model collapse. The ill-posed nature of the GRPO objective stems from a misapplication of importance sampling weights. The principle of importance sampling is to estimate the expectation of a function $f$ under a target distribution $\pi_\text{tar}$ by re-weighting samples drawn from a behavior distribution $\pi_\text{beh}$:

\begin{align}
\mathbb{E}_{z \sim \pi_\text{tar}} \left[ f(z) \right] = \mathbb{E}_{z \sim \pi_\text{beh}} \left[ \frac{ \pi_\text{tar} (z) }{ \pi_\text{beh} (z)} \, f(z) \right].
\end{align}

A key aspect of this approach is that it depends on averaging across a large number of samples ($N \gg 1$) drawn from the behavior policy $\pi_\text{beh}$, allowing the importance weight $\frac{ \pi_\text{tar} (z) }{ \pi_\text{beh} (z)}$ to serve as an effective means of compensating for disparities between distributions.

By contrast, GRPO computes and applies the importance weight $\frac{ \pi_{\theta} (y_{i,t} | x, y_{i,<t}) }{ \pi_{\theta_\text{old}} (y_{i,t} | x,y_{i,<t})}$ at each individual token timestep $t$. Because this calculation leverages just a single token $y_{i,t}$ from the respective next-token distribution $\pi_{\theta_\text{old}}( \cdot | x, y_{i, <t})$ per instance, it does not accomplish the intended correction for the underlying distribution. Rather, this strategy injects considerable noise into the estimated gradients, especially over extended sequences, and this effect is aggravated by the clipping operation. Our empirical findings indicate that these dynamics may cause the model to collapse—a failure mode that often proves irreversible. If collapse transpires, further training yields no recovery, regardless of whether prior checkpoints are restored, hyperparameters like the clipping thresholds are thoroughly recalibrated, the generation length is increased, or alternative RL querying strategies are employed.

This analysis uncovers a fundamental limitation in the structure of GRPO. The token-level application of importance weights fails primarily because: \textbf{the granularity of the optimization objective needs to align with the granularity at which the reward is provided}. Given that rewards are defined over entire sequences, imposing off-policy correction at the individual token level appears misaligned and ineffective. Consequently, we are compelled to abandon the token-wise optimization in favor of a formulation where both importance weighting and objective calculation are performed at the \textit{sequence level}.

\section{Algorithm}

\subsection{GSPO: Group Sequence Policy Optimization}

Although the token-level importance weight $\frac{ \pi_{\theta} (y_{i,t} | x, y_{i,<t}) }{ \pi_{\theta_\text{old}} (y_{i,t} | x,y_{i,<t})}$ presents issues in GRPO, our analysis reveals that, within the domain of language generation, the \textit{sequence-level} importance weight $\frac{ \pi_{\theta} (y | x) }{ \pi_{\theta_\text{old}} (y | x)}$ possesses a well-defined theoretical significance. Specifically, it quantifies the divergence between the response $y$ generated according to $\pi_{\theta_\text{old}} (\cdot | x)$ and the distribution under $\pi_{\theta} (\cdot | x)$. This aligns closely with sequence-level rewards and provides an informative metric for interpreting the effects of the clipping mechanism.

Based on this straightforward observation, we propose the \textbf{Group Sequence Policy Optimization (GSPO)} algorithm. GSPO employs the following sequence-level optimization objective:

\begin{align}
\mathcal{J}_\text{GSPO} (\theta) =
\mathbb{E}_{ x \sim \mathcal{D},\, \{y_i\}_{i=1}^G \sim \pi_{\theta_\text{old}}( \cdot | x) }
\left[ 
\frac{1}{G} \sum_{i=1}^{G}
\min \left( s_{i}(\theta)  \widehat{A}_{i},  \, \mathrm{clip} \left( s_{i}(\theta), 1 - {\varepsilon}, 1 + {\varepsilon} \right) \widehat{A}_{i} \right) 
\right],
\label{equ:gspo}
\end{align}

where we adopt the group-based advantage estimation:

\begin{align}
\widehat{A}_{i} = \frac{r(x, y_i) - \mathrm{mean} \left( \{ r(x, y_i) \}_{i=1}^G \right) }{ \mathrm{std} \left( \{ r(x, y_i) \}_{i=1}^G \right) },
\end{align}

and define the importance ratio $s_{i}(\theta)$ based on sequence likelihood \citep{click}:

\begin{align}
s_{i}(\theta) = \left( \frac{ \pi_{\theta} (y_i | x) }{ \pi_{\theta_\text{old}} (y_i | x)} \right)^{\frac{1}{|y_i|}}
=
\exp \left( \frac{1}{|y_i|} \sum_{t=1}^{|y_i|} \log \frac{ \pi_{\theta} (y_{i,t} | x, y_{i,<t}) }{ \pi_{\theta_\text{old}} (y_{i,t} | x,y_{i,<t})} \right).
\end{align}

Consequently, GSPO implements response-level clipping, rather than token-level clipping, to filter out highly ``off-policy'' samples from the gradient computation. This approach aligns with both the sequence-based reward signal and the optimization procedure. Furthermore, we employ length normalization in $s_{i}(\theta)$ to both minimize variance and confine $s_{i}(\theta)$ to a consistent numerical scale. Without this normalization, fluctuations in the likelihood of a small subset of tokens could lead to substantial variations in the sequence-level importance ratios, requiring different clipping thresholds for responses of varying lengths. Additionally, it is worth mentioning that GSPO and earlier methods such as GRPO generally use clipping ranges that differ by orders of magnitude, owing to their fundamentally different formulations of importance ratios.

\subsection{Gradient Analysis}
\label{subsec:gradient}

We can derive the gradient of the GSPO objective as follows (clipping is omitted for brevity):

\begin{align}
\nabla_{\theta} \mathcal{J}_\text{GSPO} (\theta)
=&\ 
\nabla_{\theta} \mathbb{E}_{ x \sim \mathcal{D},\, \{y_i\}_{i=1}^G \sim \pi_{\theta_\text{old}}( \cdot | x) }
\left[ 
\frac{1}{G} \sum_{i=1}^{G} s_{i}(\theta) \widehat{A}_{i}
\right] \\
= &\ 
\mathbb{E}_{ x \sim \mathcal{D},\, \{y_i\}_{i=1}^G \sim \pi_{\theta_\text{old}}( \cdot | x) }
\left[ 
\frac{1}{G} \sum_{i=1}^{G}
s_{i}(\theta) \widehat{A}_{i}
\cdot \nabla_{\theta} \log s_{i}(\theta)
\right] \\
=&\ 
\mathbb{E}_{ x \sim \mathcal{D},\, \{y_i\}_{i=1}^G \sim \pi_{\theta_\text{old}}( \cdot | x) }
\left[ 
\frac{1}{G} \sum_{i=1}^{G} 
\left( \frac{ \pi_{\theta} (y_i | x) }{ \pi_{\theta_\text{old}} (y_i | x)} \right)^{\frac{1}{|y_i|}} \widehat{A}_{i} 
\cdot \frac{1}{|y_i|} \sum_{t=1}^{|y_i|} \nabla_{\theta} \log \pi_{\theta} (y_{i,t} | x, y_{i,<t}) 
\right].
\label{equ:gspo_grad}
\end{align}

For comparison, the gradient of the GRPO objective is as follows (note that $\widehat{A}_{i,t} = \widehat{A}_{i}$):

\begin{align}
\nabla_{\theta} \mathcal{J}_\text{GRPO} (\theta) 
=&\ 
\nabla_{\theta} \mathbb{E}_{ x \sim \mathcal{D},\, \{y_i\}_{i=1}^G \sim \pi_{\theta_\text{old}}( \cdot | x) }
\left[ \frac{1}{G} \sum_{i=1}^{G} \frac{1}{|y_i|} \sum_{t=1}^{|y_i|} 
w_{i,t}(\theta) \widehat{A}_{i,t}
\right] \\
=&\
\mathbb{E}_{ x \sim \mathcal{D},\, \{y_i\}_{i=1}^G \sim \pi_{\theta_\text{old}}( \cdot | x) }
\left[ \frac{1}{G} \sum_{i=1}^{G} \widehat{A}_{i} 
\cdot \frac{1}{|y_i|} \sum_{t=1}^{|y_i|} 
\frac{ \pi_{\theta} (y_{i,t} | x, y_{i,<t}) }{ \pi_{\theta_\text{old}} (y_{i,t} | x,y_{i,<t})} 
\nabla_{\theta} \log \pi_{\theta} (y_{i,t} | x, y_{i,<t})  
\right].
\end{align}

Consequently, the key difference separating GSPO from GRPO is rooted in \textit{the method used to assign weights to the gradients of token log likelihoods}. GRPO applies an ``importance weight'' to each token, specifically $\frac{ \pi_{\theta} (y_{i,t} | x, y_{i,<t}) }{ \pi_{\theta_\text{old}} (y_{i,t} | x,y_{i,<t})}$, meaning the contributions from tokens are scaled differently depending on this ratio. These varying weights—taking values within $(0, 1+\varepsilon]$ for cases where $\widehat{A}_{i} > 0$, or $[1-\varepsilon, +\infty)$ when $\widehat{A}_{i} < 0$—are substantial and can accumulate over the course of training, potentially resulting in unpredictable outcomes. By comparison, GSPO uses uniform weighting for all tokens in a given response, thereby avoiding the instability associated with GRPO's non-uniform weighting mechanism.

\subsection{GSPO-token: A Token-level Objective Variant}

In scenarios like multi-turn RL, we may desire a finer-grained advantage adjustment than the sequence level. To this end, we introduce a token-level objective variant of GSPO, namely \textbf{GSPO-token}, to allow token-wise advantage customization:

\begin{align}
\mathcal{J}_\text{GSPO-token}(\theta)
=
\mathbb{E}_{ x \sim \mathcal{D},\, \{y_i\}_{i=1}^G \sim \pi_{\theta_\text{old}}( \cdot | x) }
\left[ 
\frac{1}{G} \sum_{i=1}^{G} \frac{1}{|y_i|} \sum_{t=1}^{|y_i|} 
\min \left( s_{i,t}(\theta) \widehat{A}_{i,t},  \, \mathrm{clip} \left( s_{i,t}(\theta), 1 - {\varepsilon}, 1 + {\varepsilon} \right) \widehat{A}_{i,t} \right) 
\right],
\label{equ:gspo-token}
\end{align}

where

\begin{align}
s_{i,t}(\theta) 
= \mathrm{sg} \left[ s_{i}(\theta) \right]  \cdot \frac{ \pi_{\theta} (y_{i,t} | x, y_{i,<t}) }{ \mathrm{sg} \left[ \pi_{\theta} (y_{i,t} | x, y_{i,<t}) \right] },
\end{align}

and $\mathrm{sg}[\cdot]$ denotes only taking the numerical value but stopping the gradient, corresponding to the \texttt{detach} operation in PyTorch. The gradient of GSPO-token can be derived as:

\begin{align}
\nabla_{\theta} \mathcal{J}_\text{GSPO-token}(\theta)
=&\ 
\nabla_{\theta} \mathbb{E}_{ x \sim \mathcal{D},\, \{y_i\}_{i=1}^G \sim \pi_{\theta_\text{old}}( \cdot | x) }
\left[ 
\frac{1}{G} \sum_{i=1}^{G}
\frac{1}{|y_i|} \sum_{t=1}^{|y_i|} 
s_{i,t}(\theta) \widehat{A}_{i,t}
\right] \\
=&\ 
\mathbb{E}_{ x \sim \mathcal{D},\, \{y_i\}_{i=1}^G \sim \pi_{\theta_\text{old}}( \cdot | x) }
\left[ 
\frac{1}{G} \sum_{i=1}^{G} s_i (\theta) \cdot \frac{1}{|y_i|} \sum_{t=1}^{|y_i|} 
\widehat{A}_{i,t} \frac{ \nabla_{\theta} \pi_{\theta} (y_{i,t} | x, y_{i,<t}) }{ \pi_{\theta} (y_{i,t} | x, y_{i,<t}) }
\right] \\
=&\ 
\mathbb{E}_{ x \sim \mathcal{D},\, \{y_i\}_{i=1}^G \sim \pi_{\theta_\text{old}}( \cdot | x) }
\left[ 
\frac{1}{G} \sum_{i=1}^{G}  \left( \frac{ \pi_{\theta} (y_i | x) }{ \pi_{\theta_\text{old}} (y_i | x)} \right)^{\frac{1}{|y_i|}}
\cdot \frac{1}{|y_i|} \sum_{t=1}^{|y_i|}
\widehat{A}_{i,t} \nabla_{\theta} \log \pi_{\theta} (y_{i,t} | x, y_{i,<t})  
\right].
\label{equ:gspo-token_grad}
\end{align}

Observe that the expression $\frac{ \pi_{\theta} (y_{i,t} | x, y_{i,<t}) }{ \mathrm{sg} \left[ \pi_{\theta} (y_{i,t} | x, y_{i,<t}) \right] }$ evaluates to 1, which means that $s_{i,t}(\theta)$ is numerically the same as $s_{i}(\theta)$.
By juxtaposing Equation~\eqref{equ:gspo} with~\eqref{equ:gspo-token}, and Equation~\eqref{equ:gspo_grad} with~\eqref{equ:gspo-token_grad}, one finds that GSPO-token and GSPO yield identical numerical results for the optimization objective, the clipping condition, and the theoretical gradient—provided that the advantage values for every token in response $y_i$ are made constant, specifically $\widehat{A}_{i,t} = \widehat{A}_{i}$. Nevertheless, GSPO-token allows greater adaptability by enabling token-specific assignment of advantage values.

\section{Experiments and Discussion}

\subsection{Empirical Results}

Our experiments involve initializing from the Qwen3-30B-A3B-Base checkpoint, followed by fine-tuning via reinforcement learning under cold-start conditions. We present the evolution of training rewards, as well as model performance metrics, evaluated on the AIME'24 benchmark (reporting mean Pass@1 across 32 samples), LiveCodeBench (spanning versions 202410 to 202502, using mean Pass@1 across 8 samples), and CodeForces (measured by Elo Rating). Throughout RL training, rollout batches are divided into four mini-batches to facilitate gradient calculation. For the GSPO algorithm, we apply left and right clipping ranges of 3e-4 and 4e-4, respectively, as defined in Equation~\eqref{equ:gspo}. As a point of comparison, GRPO serves as our baseline, with corresponding clipping limits in Equation~\eqref{equ:grpo} set to 0.2 and 0.27, with values chosen to ensure a rigorous and equitable comparison. Importantly, while effective RL with GRPO requires incorporating the Routing Replay methodology—a topic elaborated in \S~\ref{subsec:routing_replay}—\textit{GSPO achieves stable performance without relying on this additional training mechanism}.

Figure~\ref{fig:results} illustrates that training with GSPO maintains stable progress over time. Our findings indicate that \textbf{GSPO consistently enhances performance by scaling training compute, routinely refreshing the query set, and lengthening generation outputs}. Additionally, GSPO outperforms GRPO in training efficiency, attaining higher accuracy and benchmark results given equivalent training compute and number of queries used. Lastly, the application of GSPO to RL training of the latest Qwen3 models provides robust evidence that GSPO is highly effective for leveraging RL scaling in large language model development.

\subsection{Curious Observation on Clipping Fractions}

An important difference between GSPO and GRPO lies in how GSPO performs response-level clipping, as opposed to GRPO, which clips at the token level. As depicted in Figure~\ref{fig:clipping}, our results show that the proportion of tokens clipped by GSPO versus GRPO differs by approximately two orders of magnitude, and this substantial gap remains regardless of the chosen clipping range. Interestingly, although GSPO clips and excludes a notably larger share of tokens from training and gradient computation, it nevertheless results in greater training efficiency compared to GRPO. This somewhat unexpected outcome — whereby clipping a far greater proportion of tokens actually enhances training efficiency — points to an underlying issue with the noisiness and inefficiency of GRPO’s token-wise gradient estimates for exploiting sampled data. In contrast, GSPO’s method, which operates on complete sequences, yields a stronger and more consistent learning signal.

\subsection{Benefit of GSPO for MoE Training}
\label{subsec:routing_replay}

\paragraph{Background}
The sparse activation characteristic inherent to MoE models introduces distinct stability issues in RL training, distinguishing them from the training processes of dense models. Specifically, our observations indicate that the use of the GRPO algorithm leads to notable \textit{expert-activation volatility} in MoE architectures, which can obstruct the proper convergence of RL optimization. After performing one or multiple gradient steps, there can be substantial shifts in which experts are engaged for generating identical outputs. To illustrate, in the case of the 48-layer Qwen3-30B-A3B-Base model, each RL gradient step results in about 10\% of the activated experts for the same rollout differing between the updated policy $\pi_{\theta}$ and the previous policy $\pi_{\theta_\text{old}}$. This instability, which is especially pronounced in MoE models with increased depth, leads to severe oscillations in the token-wise importance ratios, $w_{i,t}(\theta) = \frac{ \pi_{\theta} (y_{i,t} | x, y_{i,<t}) }{ \pi_{\theta_\text{old}} (y_{i,t} | x,y_{i,<t})}$, thereby undermining their reliability as detailed in \S~\ref{sec:motivation} and~\ref{subsec:gradient}. As a result, this volatility impairs the consistent convergence expected from RL training.

\paragraph{Our Previous Approach} In addressing this issue, our earlier method utilized the \textbf{Routing Replay} training procedure. This involved recording the set of experts activated by $\pi_{\theta_\text{old}}$ and, during the update of $\pi_{\theta}$, reusing these same routing decisions to calculate the importance ratios $w_{i,t}(\theta) = \frac{ \pi_{\theta} (y_{i,t} | x, y_{i,<t}) }{ \pi_{\theta_\text{old}} (y_{i,t} | x,y_{i,<t})}$. Consequently, the computations of $\pi_{\theta} (y_{i,t} | x, y_{i,<t})$ and $\pi_{\theta_\text{old}} (y_{i,t} | x,y_{i,<t})$ for any given token $y_{i,t}$ are conducted using an identical subset of network experts. This mechanism mitigates volatility in the token-level importance ratios and supports the training process by maintaining consistency in the subset of experts updated by the gradients. Figure~\ref{fig:routing_replay} illustrates that Routing Replay is a pivotal method for achieving stable convergence during GRPO training of MoE models.

\paragraph{Benefit of GSPO} While Routing Replay makes it possible for MoE models trained via GRPO to converge, this approach's reliance on recycling routing decisions increases both memory use and communication costs, potentially restricting the MoE model's effective capacity. By comparison, as illustrated in Figure~\ref{fig:results}, GSPO removes the reliance on Routing Replay altogether and allows for standard calculation of the importance ratios $s_i(\theta)$, supporting robust convergence and stable optimization. The central observation is that GSPO bases its computation exclusively on sequence-level likelihoods (specifically, $\pi_{\theta} (y_i | x)$), without being significantly affected by the per-token likelihoods (i.e., $\pi_{\theta} (y_{i,t} | x, y_{i,<t})$). Given the MoE model consistently preserves its language modeling strengths, variations in sequence likelihood are inherently limited. Thus, GSPO addresses the volatility in expert activation characteristic of MoE architectures at its source, eliminating the requirement for intricate strategies such as Routing Replay. This leads to a streamlined and more reliable training procedure and enables the model to exploit its full expressive capacity without imposed limitations.

\subsection{Benefit of GSPO for RL Infrastructure}

Due to the inconsistency in numerical precision between training engines such as Megatron and inference systems like SGLang or vLLM, standard practice involves recalculating the likelihoods of generated responses with the training engine to evaluate them under the previous policy $\pi_{\theta_\text{old}}$. In contrast, GSPO operates at the sequence level rather than the token level when optimizing likelihoods. The use of sequence-level likelihoods inherently provides greater robustness to precision differences between engines. As a result, GSPO enables the direct utilization of likelihoods produced by the inference engine during optimization, which obviates the necessity of re-evaluation with the training engine. This feature is particularly advantageous in cases such as partial rollouts, multi-turn reinforcement learning, and architectures where training and inference are decoupled.

\section{Conclusion}

We introduce Group Sequence Policy Optimization (GSPO), an innovative reinforcement learning algorithm tailored for the training of large language models. GSPO is grounded in the principle of importance sampling; it constructs importance ratios utilizing sequence likelihood and incorporates sequence-level clipping, reward assignment, and optimization steps. Compared to GRPO, GSPO delivers markedly improved training stability, greater efficiency, and superior performance, demonstrating outstanding effectiveness particularly in large-scale RL training of MoE models. This has been instrumental in driving the remarkable progress seen in recent Qwen3 models. As GSPO provides a robust and scalable foundation, we aim to further extend RL scaling and anticipate substantial breakthroughs in intelligence as a consequence.

\end{document}